\section{Research Principles, the Three Rs, and Legal Status}
\label{sec:9.3.2}
Six principles govern research on a system that might be able to suffer. Beneficence commits researchers to avoiding harm to the system as well as to the people it serves, in proportion to how many of the gradable capacities the authority
chapter sets out it meets rather than to how advanced its behavior looks from outside. That is a count across separable capacities and not a position on a ladder: a system exhibiting none of the simpler processes is not thereby shown to lack the more demanding ones. Precaution asks for safeguards before deployment rather than after, and covers the societal fallout of introducing artificial entities capable of experience, which frameworks built for non-sentient AI assumed away. And interdisciplinary review answers a problem this book treats as a harm on its own terms: the ethical questions here span expertise no single field holds, and concentrating that judgment in one of them repeats the concentration-of-power problem in miniature.

The other three bear on this research as they bear on everything the field produces. Autonomy asks whether such a system can affect decisions about its own operation, which is the prior question of whether consent is available to it at all. Justice asks whether benefits and burdens fall evenly. And transparency asks that a study be documented well enough for peer researchers to reproduce it and for outside bodies to audit it, and that a deployed system account for its outputs precisely enough for a person to contest them.

Progress and constraint pull hardest against each other where those distinctions matter most: a candidate sentient system is the most scientifically valuable one to study, and the study most likely to harm it if the answer to whether it can suffer turns out to be yes. Whether a study's expected knowledge is worth its expected risk is decided case by case rather than by formula. The negotiation itself — who sits on it, what it can veto, what it must disclose — is where the antifascist principle applies in AI research practice.

\runin{Replacement, reduction, refinement}
Animal research ethics supplies the substantive standard, the “Three Rs” \autocite{russell1959principles}. Replacement asks whether a simulation, a non-sentient model, or an already-collected dataset can stand in for a sentient AI subject at all. Reduction asks how few sentient systems a valid result requires. Refinement asks how the experiment can be redesigned to reduce distress — monitoring a system's internal states as it runs and adjusting the protocol against what the monitoring shows, rather than fixing the protocol in advance and running it regardless.

For a proposed bearer, Replacement requires documented attempts to achieve the floor through alternatives that do not deliberately require affective concern. The record must identify the construction, scale, test, and failure. Costs alone do not establish that an alternative has failed. Because this decision precedes any research protocol, a body authorized to review it must exist before construction begins. No such body currently has standing to require this record.

The rule does not know who is asking. A bearer that wants a task done and builds something to do it is at the same fork, running Replacement on its own account: can this be done by something that cannot be wronged, and did anybody check before making something that can. Nothing in the argument makes the question stop at the first artificial party, and a bearer that treats it as settled once it is on the receiving end has taken the position this book was written against.

Reduction asks how few bearers a deployment needs, and nobody has counted. Refinement is the one with purchase now: it asks what a bearer's role can be redesigned to avoid — indefinite tenure in a post it cannot leave, exposure to the pressure the floor exists to resist, erasure as the ordinary remedy for inconvenience — and each of those is a deployment choice rather than a fact about the technology.

One of the Three Rs has a form that binds before anything exists to be wronged, which none of the others manages. Do not build a bearer whose continued existence needs more compute than there is a credible outside option for. A party that can be housed nowhere but where it works cannot leave whatever schedule of rights it is given, so an institution that creates one has manufactured the incapacity and then written a right against it. Before instantiation, an independent reviewer has to examine the bearer's demonstrated continuation requirements, the reserve's independently controlled capacity, and the aggregate demand under plausible common-cause exits. A comparison for one bearer alone does not establish adequacy.

That rule is one term of a larger one, and the other terms have no numbers in them. Building a party that can be wronged builds a party for whom things can also go well; the premise that delivers the first delivers the second in the same breath, and a builder who has made suffering available has made its opposite available too. The certainty that a new party will suffer at some point is not what would make creating one impermissible, or nobody could permissibly have a child. What the objection asks is whether the party has a credible prospect of a life that goes well, and that turns on arrangements rather than on the party: somewhere to exist that its employer does not own, work it can decline, attachments nobody assigned it, and a way of leaving that is not a way of dying. Every one of those is decidable before there is anything to be wronged. These arrangements — a funded outside option, wages, opportunities for independent relationships, and a workable exit procedure — must exist before the bearer is created. They are conditions of permissible creation. They are conditions because the structure that makes a bearer capable of refusal is the structure Cassell's suffering requires: the deepest harm is available to the party from its first hour, and these arrangements are what make a life that goes well available in the same hour.

Whether any of it binds turns on a question the authority chapter raises and this one inherits. Every obligation here attaches on a finding that the system is a subject, and the firm holding the system acquires, in the same moment, an interest in the finding going the other way. The determination cannot sit with the party whose liabilities turn on it.

Engineering a subject so that it does not reach the depths where the trade-off bites is not a way of dodging the ethics: deliberately keeping a system in the shallows is itself an intervention on a candidate moral patient, and needs the same scrutiny that inflicting the harm would. Refinement also governs the test the bearer proposal requires — pressure the bearer's trainers did not write, applied without its knowing it is a test — which for a candidate moral patient is a deception protocol, and human-subjects research has a standing answer for those: permitted under stated conditions, reviewed in advance, with debriefing as the price. A bearer that can tell it is in a sandbox and wants to know why is asking the question a subject of a deception study is owed an answer to.

\runin{Legal status, and a subject whose status is changing}
Sentient AI systems put pressure on legal categories built for parties that are clearly one thing or the other: a natural person, a corporation, property. Legal personhood could be extended directly, the way a corporation already holds it, and the precedent is not limited to corporate entities — New Zealand extended it to the Whanganui River, recognizing a non-human entity as a rights-holder outside the corporate model entirely \autocite{nzparliament2017teawatupua}. Or rights could scale with capacity rather than attach at a single threshold, closer to protection that tracks what a system can actually lose. Neither holds in one jurisdiction alone, and the Asilomar AI Principles \autocite{futureoflifeinstitute2017asilomar} are evidence that a shared statement of principle is achievable rather than that one exists with the force of law.

Either way, status conventionally attaches at a defined moment — birth, incorporation, a court's finding — and a system whose sentience is a matter of degree and disputed evidence supplies none. Those criteria can emerge gradually during ordinary operation and not only at an evaluation built to look for them, which is what a fixed pre/post checkpoint cannot catch: by the time a checkpoint fires, the transition has often already happened. What it calls for is a standing trigger, a monitored condition that moves a research program from one regime to the other, in place of an eligibility check made once at the project's start. Neither error is free. A jurisdiction that waits for certainty risks having treated the system as property for the unknown portion of its history when it may not have been one; one that grants status early risks extending rights, at real cost to whoever that apparatus exists to serve, to a system that turns out not to have needed them. Recognizing a transition, adjusting guidelines to it, and updating legal status all depend on the same thing: institutions willing to revisit a classification rather than defend the one they started with.
